\chapter{All-Pairs Shortest Paths}

In the previous chapter you learned how to stand at one fixed vertex -- call it
the \emph{source} -- and compute the shortest distance from that single source
to every other vertex in a graph. Dijkstra's algorithm did this when all the
edge weights were nonnegative; the Bellman--Ford algorithm did it even when
some weights were negative, and it had the good manners to tell you when a
negative-weight cycle made the whole question meaningless.

But sometimes one source is not enough. A road atlas in the back of an old
travel guide does not give you distances from a single city; it gives you a
big triangular table where you can look up the distance between \emph{any} two
cities at once. A router wants to know the best next hop toward every possible
destination, not just one. A program that scores how ``central'' each person is
in a social network needs the distance between every pair of people. In all of
these you want the shortest path between \emph{every} ordered pair of vertices.
That is the all-pairs shortest-paths problem, and it is the subject of this
chapter.

\begin{intuitionbox}
You already own a hammer that solves this: run a single-source algorithm once
from each vertex. With $n$ vertices, run Dijkstra $n$ times, or Bellman--Ford
$n$ times. That genuinely works, and for some graphs it is the best thing to
do. But running $n$ separate searches feels wasteful -- each run rediscovers
structure the others also needed. The fresh idea of this chapter is to forget
about ``search from a source'' entirely and instead think about the answer as
one big table that we refine, all $n^2$ entries at once, in a handful of
sweeps. That table-refinement view turns out to be pure dynamic programming,
and it produces two of the most quietly elegant algorithms in the book:
the matrix-multiplication method and Floyd--Warshall.
\end{intuitionbox}

\begin{focusbox}
Hold on to three numbers as you read. The input is a graph with $n$ vertices
and $m$ edges. The output is an $n \times n$ matrix $D$ where $D[i][j]$ is the
weight of a shortest path from $i$ to $j$. The repeated-Dijkstra baseline costs
$\Oh(n m \log n)$ (or $\Oh(n m + n^2 \log n)$ with a Fibonacci heap);
Floyd--Warshall costs a clean $\Th(n^3)$ no matter what; and Johnson's
algorithm gets you down to roughly $\Oh(n m \log n)$ \emph{even with negative
edges}, which is the best of both worlds for sparse graphs.
\end{focusbox}

\section{Setting Up: Inputs, Outputs, and Conventions}
\label{ch25:setup}

Let us fix notation once and reuse it for the whole chapter. We are given a
directed, weighted graph $G = (V, E)$ with $V = \{1, 2, \dots, n\}$. Each edge
$(i,j) \in E$ carries a real weight $w(i,j)$, which may be negative. We will
represent the graph as an $n \times n$ \emph{weight matrix} $W$, defined by

\[
W[i][j] =
\begin{cases}
0 & \text{if } i = j,\\[2pt]
w(i,j) & \text{if } i \neq j \text{ and } (i,j) \in E,\\[2pt]
\infty & \text{if } i \neq j \text{ and } (i,j) \notin E.
\end{cases}
\]

The diagonal is $0$ because the empty path from a vertex to itself has weight
zero (assuming no negative-weight cycles, which we discuss below). A missing
edge is $\infty$ because you cannot get there in one hop. This matrix is the
raw material; everything else is built from it.

The output is the \emph{distance matrix} $D$, where $D[i][j] = \delta(i,j)$,
the weight of a shortest path from $i$ to $j$, with the usual convention that
$\delta(i,j) = \infty$ if $j$ is unreachable from $i$.

\begin{trapbox}
A negative-weight \emph{cycle} reachable on a path from $i$ to $j$ makes
$\delta(i,j) = -\infty$: you can loop the cycle forever, shaving off weight
each time, so there is no shortest path at all. Negative edges by themselves
are fine; negative \emph{cycles} are the poison. Every algorithm in this
chapter assumes there is no negative-weight cycle, and we will note how each
one can be made to \emph{detect} one if it is present.
\end{trapbox}

We will also want a second output: not just the distances, but the actual
paths. For that we keep a \emph{predecessor matrix} $\Pi$, where $\Pi[i][j]$ is
the vertex that comes just before $j$ on a shortest path from $i$ to $j$ (or
\textsc{nil} if no path exists or if $i = j$). Reading off the path then means
walking backwards: from $j$ to $\Pi[i][j]$ to $\Pi[i][\Pi[i][j]]$ and so on
until you reach $i$. The $i$-th row of $\Pi$ is exactly a shortest-paths tree
rooted at $i$, the same object Dijkstra and Bellman--Ford produced.

\begin{intuitionbox}
Think of the whole chapter as bookkeeping with two matrices that grow up
together. $D$ holds the best distances we currently believe; $\Pi$ holds the
``how did I get here?'' breadcrumbs. Every time we improve a $D$ entry we also
update its $\Pi$ entry to remember the move that made the improvement. If you
ever lose the thread, ask: \emph{what does $D[i][j]$ mean right now, and what
does $\Pi[i][j]$ point at?}
\end{intuitionbox}

\section{First Idea: Shortest Paths as Matrix ``Multiplication''}
\label{ch25:matrix}

We will warm up with an algorithm that is a little slower than what we
ultimately want, but that reveals the dynamic-programming skeleton with
unusual clarity. The trick is to organize shortest paths by \emph{how many
edges they use}.

\subsection{The subproblem: shortest paths with a bounded number of edges}

Let $\ell^{(r)}_{ij}$ denote the minimum weight of any path from $i$ to $j$
that uses \emph{at most} $r$ edges. This is our subproblem. The base case is
the zero-edge case:

\[
\ell^{(0)}_{ij} =
\begin{cases}
0 & \text{if } i = j,\\
\infty & \text{if } i \neq j.
\end{cases}
\]

With zero edges you can only stay put. Now the recursive step. Suppose we
already know $\ell^{(r-1)}_{ij}$ for all pairs -- the best you can do with at
most $r-1$ edges. How do we get the best path with at most $r$ edges from $i$
to $j$? Either we are happy with a path of at most $r-1$ edges (which is
already covered), or we take a path of at most $r-1$ edges from $i$ to some
intermediate vertex $k$ and then add the single final edge $(k, j)$. We try
every possible last-but-one vertex $k$ and keep the best:

\[
\ell^{(r)}_{ij}
= \min\left(\ell^{(r-1)}_{ij},\ \min_{1 \le k \le n}\ \big(\ell^{(r-1)}_{ik} + W[k][j]\big)\right).
\]

Here is the observation that gives this section its name. Because
$W[j][j] = 0$, the term $k = j$ inside the inner $\min$ reproduces
$\ell^{(r-1)}_{ij}$ exactly. So the outer $\min$ with $\ell^{(r-1)}_{ij}$ is
redundant, and we can write the whole thing as one tidy minimization:

\[
\ell^{(r)}_{ij} = \min_{1 \le k \le n}\ \big(\ell^{(r-1)}_{ik} + W[k][j]\big).
\]

Stare at that formula and compare it with ordinary matrix multiplication,
$C[i][j] = \sum_{k} A[i][k] \cdot B[k][j]$. They are the same shape! We have
just replaced ``$+$'' with ``$\min$'' and ``$\times$'' with ``$+$''. In this
\emph{min-plus} algebra (also called the tropical semiring), computing the
next layer of distances \emph{is} a matrix product. If we name the operation
$\odot$, then $L^{(r)} = L^{(r-1)} \odot W$.

\begin{formalbox}
Define the min-plus product of two $n \times n$ matrices $A$ and $B$ by
\[
(A \odot B)[i][j] = \min_{1 \le k \le n}\ \big(A[i][k] + B[k][j]\big).
\]
Let $L^{(1)} = W$ (at most one edge) and $L^{(r)} = L^{(r-1)} \odot W$. Then
$L^{(r)}[i][j] = \ell^{(r)}_{ij}$, the minimum weight of an $i$-to-$j$ path
using at most $r$ edges.
\end{formalbox}

Figure~\ref{ch25:fig-minplus} shows the shape of one entry of the product: to
fill $L'[i][j]$ we slide across row $i$ of $L$ and down column $j$ of $W$, add
the paired entries, and keep the smallest sum -- the same dance as an ordinary
dot product, but with $\min$ for the outer combine and $+$ for the inner one.

\begin{figure}[H]
\centering
\begin{tikzpicture}[>=Stealth, scale=0.92, every node/.style={transform shape}]
  % L matrix (row i highlighted)
  \matrix (Lm) [matrix of math nodes, nodes={minimum size=6mm}, left delimiter=(, right delimiter=)] {
    \cdot & \cdot & \cdot \\
    L[i][1] & L[i][2] & L[i][3] \\
    \cdot & \cdot & \cdot \\
  };
  \node[above=1mm of Lm] {\small $L$ (row $i$)};
  \begin{scope}[on background layer]
    \node[fill=sourceblue!18, fit=(Lm-2-1)(Lm-2-3)] {};
  \end{scope}
  % odot symbol
  \node[right=4mm of Lm] (op) {$\odot$};
  % W matrix (column j highlighted)
  \matrix (Wm) [right=4mm of op, matrix of math nodes, nodes={minimum size=6mm}, left delimiter=(, right delimiter=)] {
    \cdot & W[1][j] & \cdot \\
    \cdot & W[2][j] & \cdot \\
    \cdot & W[3][j] & \cdot \\
  };
  \node[above=1mm of Wm] {\small $W$ (col $j$)};
  \begin{scope}[on background layer]
    \node[fill=targetgreen!18, fit=(Wm-1-2)(Wm-3-2)] {};
  \end{scope}
  % result
  \node[below=6mm of op] (eq) {$\Rightarrow$};
  \node[below=2mm of eq, draw, fill=quotientorange!18, rounded corners] (res)
    {$L'[i][j]=\displaystyle\min_{k}\big(L[i][k]+W[k][j]\big)$};
\end{tikzpicture}
\caption{One entry of the min-plus product $L' = L \odot W$. Slide row $i$ of
$L$ against column $j$ of $W$, add paired cells, and keep the minimum.}
\label{ch25:fig-minplus}
\end{figure}

\subsection{When do we stop?}

In a graph with no negative-weight cycles, any shortest path is \emph{simple}:
it never repeats a vertex, because cutting out a loop can only keep the weight
the same or lower it. A simple path among $n$ vertices has at most $n-1$ edges.
So once $r$ reaches $n-1$, the table stops changing: $L^{(n-1)} = L^{(n)} =
\cdots$, and $L^{(n-1)}[i][j] = \delta(i,j)$. We have the answer.

\begin{labbox}
Let us trace it on a tiny graph. Vertices $\{1,2,3\}$ with edges
$1 \to 2$ (weight $3$), $2 \to 3$ (weight $4$), and $1 \to 3$ (weight $9$).
The weight matrix (rows $=$ from, columns $=$ to; ``$\cdot$'' is $\infty$) is
\[
W = L^{(1)} =
\begin{pmatrix}
0 & 3 & 9\\
\cdot & 0 & 4\\
\cdot & \cdot & 0
\end{pmatrix}.
\]
Now compute $L^{(2)} = L^{(1)} \odot W$. The interesting entry is
$L^{(2)}[1][3] = \min_k (L^{(1)}[1][k] + W[k][3])$. Try each $k$:
$k=1$ gives $0 + 9 = 9$; $k=2$ gives $3 + 4 = 7$; $k=3$ gives $9 + 0 = 9$.
The minimum is $7$, found via $k = 2$. So routing through vertex $2$ beats the
direct edge:
\[
L^{(2)} =
\begin{pmatrix}
0 & 3 & 7\\
\cdot & 0 & 4\\
\cdot & \cdot & 0
\end{pmatrix}.
\]
With $n = 3$ we stop at $r = n-1 = 2$, and indeed $D = L^{(2)}$. The shortest
path from $1$ to $3$ has weight $7$, going $1 \to 2 \to 3$.
\end{labbox}

\subsection{Pseudocode and a first cost}

\begin{algorithm}[H]
\caption{\textsc{Extend-Shortest-Paths}($L$, $W$) -- one min-plus product}
\KwIn{An $n \times n$ matrix $L$ (best distances so far) and weight matrix $W$.}
\KwOut{The matrix $L' = L \odot W$.}
let $L'$ be a new $n \times n$ matrix\;
\For{$i \leftarrow 1$ \KwTo $n$}{
  \For{$j \leftarrow 1$ \KwTo $n$}{
    $L'[i][j] \leftarrow \infty$\;
    \For{$k \leftarrow 1$ \KwTo $n$}{
      $L'[i][j] \leftarrow \min\big(L'[i][j],\ L[i][k] + W[k][j]\big)$\;
    }
  }
}
\Return $L'$\;
\end{algorithm}

One product costs $\Th(n^3)$, exactly like grade-school matrix multiplication.
We need products $L^{(2)}$, $L^{(3)}$, all the way up to $L^{(n-1)}$, which is
$n - 2$ of them,
so the naive method costs $\Th(n^4)$. That is worse than running Bellman--Ford
$n$ times ($\Th(n^2 m)$, which is $\Th(n^4)$ on a dense graph but better on a
sparse one), so why bother? Because of one clever speed-up.

\subsection{Repeated squaring: $\Th(n^3 \log n)$}

We do not actually need every intermediate $L^{(r)}$; we only need
$L^{(n-1)}$. And once the table stabilizes at $L^{(n-1)}$, multiplying by $W$
again changes nothing. So $L^{(r)} = L^{(n-1)}$ for every $r \ge n-1$. That
means we can leap ahead by \emph{squaring}: $L^{(2r)} = L^{(r)} \odot L^{(r)}$.

\begin{intuitionbox}
This is the same trick as fast exponentiation. To compute $x^{16}$ you do not
multiply $x$ by itself sixteen times; you square four times:
$x \to x^2 \to x^4 \to x^8 \to x^{16}$. Here we square the distance matrix
instead. Each squaring doubles the number of edges our paths may use, so after
$\lceil \log_2(n-1)\rceil$ squarings we have surpassed $n-1$ edges and the
answer has converged.
\end{intuitionbox}

\begin{algorithm}[H]
\caption{\textsc{Faster-APSP}($W$) -- min-plus repeated squaring}
\KwIn{Weight matrix $W$ of an $n$-vertex graph with no negative cycle.}
\KwOut{Distance matrix $D$.}
$L \leftarrow W$\;
$r \leftarrow 1$\;
\While{$r < n - 1$}{
  $L \leftarrow \textsc{Extend-Shortest-Paths}(L, L)$\;
  $r \leftarrow 2r$\;
}
\Return $L$\;
\end{algorithm}

Each squaring is one $\Th(n^3)$ product, and we do $\lceil \log_2(n-1) \rceil$
of them, for a total of $\Th(n^3 \log n)$. Notice we feed $L$ as \emph{both}
arguments now (squaring), not $L$ and $W$.

\begin{complexitybox}
\textbf{Matrix-multiplication APSP.}
\begin{itemize}
\item Naive layer-by-layer: $\Th(n^4)$ time, $\Th(n^2)$ space.
\item Repeated squaring: $\Th(n^3 \log n)$ time, $\Th(n^2)$ space.
\item Handles negative edges; assumes no negative-weight cycle.
\item Sub-cubic ``fast'' matrix multiplication does \emph{not} directly help,
because min-plus has no subtraction, so Strassen-style cancellation tricks do
not apply. Beating $n^3$ in min-plus is a deep open-ish problem.
\end{itemize}
\end{complexitybox}

\begin{walkbox}
We argued $L^{(r)}$ stops changing at $r = n-1$ because shortest paths are
simple. Suppose the graph \emph{did} contain a negative-weight cycle. Trace
what would happen to a $D[i][j]$ for some $i, j$ that can reach and leave that
cycle as you keep squaring. Does the table ever stabilize? What does that
suggest about how to \emph{detect} a negative cycle with this method?
\end{walkbox}

\section{Floyd--Warshall: A Smarter Dynamic Program}
\label{ch25:floyd}

The matrix method organized paths by their \emph{number of edges}. Robert Floyd
and Stephen Warshall organized them by a completely different feature, and it
buys us a factor of $\log n$, landing at a clean $\Th(n^3)$ with code so short
you can write it from memory.

\subsection{The new subproblem: allowed intermediate vertices}

Fix an ordering of the vertices $1, 2, \dots, n$. Consider a path from $i$ to
$j$. Its \emph{intermediate vertices} are all the vertices on it except the two
endpoints $i$ and $j$. The new idea: build up the set of intermediate vertices
we are \emph{allowed to use}, one vertex at a time.

Define $d^{(k)}_{ij}$ to be the weight of a shortest path from $i$ to $j$ whose
intermediate vertices all lie in the set $\{1, 2, \dots, k\}$. (The endpoints
$i$ and $j$ are always allowed; the restriction is only on what we may pass
\emph{through}.) The base case uses no intermediate vertices at all:

\[
d^{(0)}_{ij} = W[i][j].
\]

With $k = 0$ the only allowed paths are direct edges, which is exactly $W$.

\subsection{The recurrence: ``do I route through vertex $k$ or not?''}

Now suppose we know $d^{(k-1)}_{ij}$ for all $i, j$ -- shortest paths allowed
to pass through $\{1, \dots, k-1\}$. We want $d^{(k)}_{ij}$, which additionally
permits passing through vertex $k$. A shortest such path either uses $k$ as an
intermediate vertex or it does not. These two cases tile every possibility:

\begin{itemize}
\item \textbf{Does not use $k$.} Then it only passes through
$\{1, \dots, k-1\}$, so its weight is $d^{(k-1)}_{ij}$.
\item \textbf{Uses $k$ exactly once.} (At most once, since the path is simple
when there are no negative cycles.) Then it splits at $k$ into a piece from $i$
to $k$ and a piece from $k$ to $j$, and \emph{neither} piece uses $k$ as an
intermediate vertex -- they only pass through $\{1, \dots, k-1\}$. So the
weight is $d^{(k-1)}_{ik} + d^{(k-1)}_{kj}$.
\end{itemize}

Take the better of the two (Figure~\ref{ch25:fig-routek}):

\[
\boxed{\;d^{(k)}_{ij} = \min\big(d^{(k-1)}_{ij},\ \ d^{(k-1)}_{ik} + d^{(k-1)}_{kj}\big).\;}
\]

After we have allowed all $n$ vertices as intermediates, no restriction
remains, so $d^{(n)}_{ij} = \delta(i,j)$. That is the whole algorithm.

\begin{intuitionbox}
Picture the sentence: ``I now permit myself to drive through city $k$. For
each pair of cities $(i,j)$, is it worth detouring through $k$?'' Routing
through $k$ is worth it exactly when the trip $i \to k$ plus the trip $k \to j$
(each using only cities I had already cleared) is shorter than my current best
$i \to j$. Sweep $k$ from $1$ to $n$, opening up one city at a time, and after
the last sweep every city is open and every distance is final.
\end{intuitionbox}

\begin{figure}[H]
\centering
\begin{tikzpicture}[>=Stealth,node distance=2.2cm]
  \node[circle,draw,fill=sourceblue!18] (i) {$i$};
  \node[circle,draw,fill=structurepurple!18,right=of i] (k) {$k$};
  \node[circle,draw,fill=targetgreen!18,right=of k] (j) {$j$};
  \draw[->,thick] (i) to[bend left=18] node[above]{\small $d^{(k-1)}_{ik}$} (k);
  \draw[->,thick] (k) to[bend left=18] node[above]{\small $d^{(k-1)}_{kj}$} (j);
  \draw[->,dashed,gray] (i) to[bend right=30] node[below]{\small $d^{(k-1)}_{ij}$ (skip $k$)} (j);
\end{tikzpicture}
\caption{The Floyd--Warshall choice for the pair $(i,j)$ at step $k$: either
keep the best path that avoids $k$ (dashed) or detour through $k$ (solid),
gluing together the best $i$-to-$k$ and $k$-to-$j$ pieces that already avoid
$k$.}
\label{ch25:fig-routek}
\end{figure}

\subsection{We only need two layers... actually one}

The recurrence reads layer $k-1$ to produce layer $k$. Naively that suggests
storing all $n+1$ layers, costing $\Th(n^3)$ space. But you can do better:
\emph{overwrite the matrix in place}. Use a single
matrix $D$ and the assignment
$D[i][j] \leftarrow \min(D[i][j],\ D[i][k] + D[k][j])$.

Why is reading possibly-already-updated values safe? Look at the two cells we
read: $D[i][k]$ and $D[k][j]$. In the row $k$ and column $k$ entries,
$d^{(k)}_{ik} = d^{(k-1)}_{ik}$ and $d^{(k)}_{kj} = d^{(k-1)}_{kj}$, because a
shortest path to or from $k$ that is allowed to use $k$ as an intermediate
gains nothing -- using $k$ in the middle would just repeat the endpoint. So the
$k$-th row and $k$-th column do not change during pass $k$, and reading them
``early'' returns the right old value. Single matrix, $\Th(n^2)$ space.

Figure~\ref{ch25:fig-fwcross} captures the picture: during pass $k$, the gray
cross (row $k$ and column $k$) is frozen, and the update to any cell $(i,j)$
reads exactly the two cross cells $D[i][k]$ and $D[k][j]$. Because those two
never move during this pass, it does not matter whether they have already been
visited.

\begin{figure}[H]
\centering
\begin{tikzpicture}[>=Stealth]
  \def\nn{5}
  % grid
  \foreach \r in {1,...,5} {
    \foreach \c in {1,...,5} {
      \node[draw, minimum size=6.5mm] (c\r-\c) at (\c*0.72, -\r*0.72) {};
    }
  }
  % freeze row k=3 and column k=3
  \begin{scope}[on background layer]
    \foreach \c in {1,...,5} { \node[fill=gray!22, minimum size=6.5mm] at (\c*0.72, -3*0.72) {}; }
    \foreach \r in {1,...,5} { \node[fill=gray!22, minimum size=6.5mm] at (3*0.72, -\r*0.72) {}; }
  \end{scope}
  % cell (i,j) = (2,5) being updated
  \node[fill=quotientorange!30, minimum size=6.5mm] at (5*0.72,-2*0.72) {$ij$};
  % the two read cells D[i][k]=(2,3) and D[k][j]=(3,5)
  \node[fill=sourceblue!30, minimum size=6.5mm] at (3*0.72,-2*0.72) {$ik$};
  \node[fill=targetgreen!30, minimum size=6.5mm] at (5*0.72,-3*0.72) {$kj$};
  % arrows
  \draw[->,thick,sourceblue] (3*0.72,-2*0.72) -- (4.6*0.72,-2*0.72);
  \draw[->,thick,targetgreen] (5*0.72,-2.6*0.72) -- (5*0.72,-2.18*0.72);
  % labels
  \node[left=1mm of c3-1] {\small row $k$};
  \node[above=0.5mm] at (3*0.72, -0.4) {\small col $k$};
\end{tikzpicture}
\caption{Pass $k$ of Floyd--Warshall on a $5\times5$ table. The frozen cross is
row $k$ and column $k$ (gray). Updating cell $(i,j)$ (orange) reads only
$D[i][k]$ (blue) and $D[k][j]$ (green), both on the frozen cross -- which is why
the in-place overwrite is safe.}
\label{ch25:fig-fwcross}
\end{figure}

\begin{algorithm}[H]
\caption{\textsc{Floyd-Warshall}($W$) with predecessor reconstruction}
\KwIn{Weight matrix $W$ of an $n$-vertex graph with no negative cycle.}
\KwOut{Distance matrix $D$ and predecessor matrix $\Pi$.}
$D \leftarrow W$\;
\For{$i \leftarrow 1$ \KwTo $n$}{
  \For{$j \leftarrow 1$ \KwTo $n$}{
    \lIf{$i \neq j$ \textbf{and} $W[i][j] < \infty$}{$\Pi[i][j] \leftarrow i$}
    \lElse{$\Pi[i][j] \leftarrow \textsc{nil}$}
  }
}
\For{$k \leftarrow 1$ \KwTo $n$}{
  \For{$i \leftarrow 1$ \KwTo $n$}{
    \For{$j \leftarrow 1$ \KwTo $n$}{
      \If{$D[i][k] + D[k][j] < D[i][j]$}{
        $D[i][j] \leftarrow D[i][k] + D[k][j]$\;
        $\Pi[i][j] \leftarrow \Pi[k][j]$\;
      }
    }
  }
}
\Return $(D, \Pi)$\;
\end{algorithm}

The predecessor update deserves a word. When routing through $k$ improves the
path from $i$ to $j$, the new path ends with the $k$-to-$j$ portion, so the
vertex just before $j$ is whatever was just before $j$ on the best $k$-to-$j$
path: that is $\Pi[k][j]$. Copying it gives correct breadcrumbs.

\begin{focusbox}
The order of the loops is sacred: $k$ \emph{outermost}, then $i$, then $j$. The
$k$ loop is the dynamic-programming dimension (which intermediates are
allowed); $i$ and $j$ just sweep all pairs within one $k$-layer. Swap $k$ to
the inside and you compute nonsense, because you would be allowing each pair a
different, inconsistent set of intermediates.
\end{focusbox}

\subsection{A full worked trace}

\begin{labbox}
Take a four-vertex directed graph with these edges (and weights), which we will
carry through:
$1\!\to\!2 = 3$, $1\!\to\!4 = 7$, $2\!\to\!1 = 8$, $2\!\to\!3 = 2$,
$3\!\to\!1 = 5$, $3\!\to\!4 = 1$, $4\!\to\!1 = 2$, $4\!\to\!3 = 6$.
The starting matrix $D = D^{(0)} = W$ (``$\cdot$'' is $\infty$):
\[
D^{(0)}=
\begin{pmatrix}
0 & 3 & \cdot & 7\\
8 & 0 & 2 & \cdot\\
5 & \cdot & 0 & 1\\
2 & \cdot & 6 & 0
\end{pmatrix}.
\]

\textbf{Pass $k=1$} (route through vertex $1$). We improve $D[i][j]$ whenever
$D[i][1] + D[1][j] < D[i][j]$. For instance $D[2][4]$: currently $\infty$;
$D[2][1] + D[1][4] = 8 + 7 = 15$, so it drops to $15$. Also $D[3][2]$:
$D[3][1] + D[1][2] = 5 + 3 = 8$, new. And $D[4][2]$: $D[4][1] + D[1][2] = 2 + 3
= 5$, new.
\[
D^{(1)}=
\begin{pmatrix}
0 & 3 & \cdot & 7\\
8 & 0 & 2 & 15\\
5 & 8 & 0 & 1\\
2 & 5 & 6 & 0
\end{pmatrix}.
\]

\textbf{Pass $k=2$} (now also allow vertex $2$). Check $D[1][3]$:
$D[1][2] + D[2][3] = 3 + 2 = 5 < \infty$, so it becomes $5$. Check $D[1][4]$:
$D[1][2] + D[2][4] = 3 + 15 = 18$, not better than $7$. $D[3][3]$ stays $0$.
\[
D^{(2)}=
\begin{pmatrix}
0 & 3 & 5 & 7\\
8 & 0 & 2 & 15\\
5 & 8 & 0 & 1\\
2 & 5 & 6 & 0
\end{pmatrix}.
\]

\textbf{Pass $k=3$} (allow vertex $3$). $D[1][4]$: $D[1][3] + D[3][4] = 5 + 1 =
6 < 7$, improves to $6$! $D[2][4]$: $D[2][3] + D[3][4] = 2 + 1 = 3 < 15$, big
improvement to $3$. $D[2][1]$: $D[2][3] + D[3][1] = 2 + 5 = 7 < 8$, drops to
$7$. $D[4][1]$ stays $2$.
\[
D^{(3)}=
\begin{pmatrix}
0 & 3 & 5 & 6\\
7 & 0 & 2 & 3\\
5 & 8 & 0 & 1\\
2 & 5 & 6 & 0
\end{pmatrix}.
\]

\textbf{Pass $k=4$} (allow vertex $4$). $D[3][2]$: $D[3][4] + D[4][2] = 1 + 5 =
6 < 8$, improves to $6$. $D[3][1]$: $D[3][4] + D[4][1] = 1 + 2 = 3 < 5$, drops
to $3$. $D[1][3]$: $D[1][4] + D[4][3] = 6 + 6 = 12$, no. $D[2][1]$:
$D[2][4]+D[4][1] = 3 + 2 = 5 < 7$, drops to $5$.
\[
D^{(4)} = D =
\begin{pmatrix}
0 & 3 & 5 & 6\\
5 & 0 & 2 & 3\\
3 & 6 & 0 & 1\\
2 & 5 & 6 & 0
\end{pmatrix}.
\]
That final matrix is the answer: every entry is a true shortest distance. For
example, the shortest way from $3$ to $1$ has weight $3$, achieved by
$3 \to 4 \to 1$ ($1 + 2$), which the direct edge ($5$) could not match.
\end{labbox}

\subsection{Why it is correct}

\begin{formalbox}
\textbf{Loop invariant (Floyd--Warshall).} Just before the iteration of the
outer loop with index $k$, every entry satisfies
\[
D[i][j] = d^{(k-1)}_{ij} = \text{(weight of a shortest } i\text{-to-}j\text{ path
through intermediates in } \{1,\dots,k-1\}).
\]
\end{formalbox}

\emph{Initialization.} Before $k = 1$ we have $D = W = D^{(0)}$: paths through
the empty set of intermediates are exactly direct edges. The invariant holds.

\emph{Maintenance.} Assume it holds entering pass $k$. The update sets
$D[i][j] \leftarrow \min(D[i][j], D[i][k] + D[k][j])$. By the in-place argument
above, the values read are the old layer-$(k-1)$ values. By the case analysis
(route through $k$ or not), this min is exactly $d^{(k)}_{ij}$. So after the
pass the invariant holds for $k+1$.

\emph{Termination.} After $k = n$ finishes, $D[i][j] = d^{(n)}_{ij}$, paths
allowed to use \emph{any} intermediate, which is the unrestricted shortest
distance $\delta(i,j)$. Done.

\begin{trapbox}
If a negative-weight cycle exists, Floyd--Warshall does not loop forever -- it
runs the same $\Th(n^3)$ steps -- but the answers become meaningless, and some
diagonal entry $D[i][i]$ ends up \emph{negative}. That is your detector: after
running, scan the diagonal; if any $D[i][i] < 0$, vertex $i$ lies on a
negative-weight cycle and the affected distances are really $-\infty$.
\end{trapbox}

\begin{complexitybox}
\textbf{Floyd--Warshall.} Three nested loops of length $n$: $\Th(n^3)$ time,
flat, best case equals worst case. Space $\Th(n^2)$ (single in-place matrix,
plus $\Pi$ if you want paths). Handles negative edges; detects negative cycles
via the diagonal. No graph is too dense for it -- in fact it is the method of
choice precisely \emph{when} the graph is dense.
\end{complexitybox}

\section{Transitive Closure: Reachability for the Price of Booleans}
\label{ch25:closure}

Sometimes you do not care about distances at all -- you only want to know
\emph{whether} a path exists from $i$ to $j$. The transitive closure of $G$ is
the graph $G^*$ with an edge $(i, j)$ exactly when $j$ is reachable from $i$ in
$G$. We adopt the convention that every vertex reaches itself (the empty path),
so the diagonal is always \texttt{true}; off the diagonal, $(i,j)$ means there
is a directed path of one or more edges. Think of ``can task $i$ eventually
trigger task $j$?'' in a dependency graph, or ``can person $i$ reach person $j$
through any chain of follows?''

You could just run Floyd--Warshall with weights $0$ and $\infty$ and then mark
``reachable'' wherever $D[i][j] < \infty$. That works. But if all you need is a
yes/no answer, you can run the \emph{same skeleton} on booleans and avoid
arithmetic entirely, which is faster in practice and uses single bits.

Define $t^{(k)}_{ij}$ to be \texttt{true} iff there is a path from $i$ to $j$
with all intermediate vertices in $\{1, \dots, k\}$. The base case is
$t^{(0)}_{ij} = \texttt{true}$ iff $i = j$ or $(i,j) \in E$. The recurrence
mirrors Floyd--Warshall, with $\min$ becoming logical OR and $+$ becoming
logical AND:

\[
t^{(k)}_{ij} = t^{(k-1)}_{ij}\ \ \textbf{or}\ \ \big(t^{(k-1)}_{ik}\ \ \textbf{and}\ \ t^{(k-1)}_{kj}\big).
\]

\begin{intuitionbox}
The translation is exact. ``Shortest distance'' lived in the
$(\min, +)$ world; ``can you get there'' lives in the $(\textsc{or},
\textsc{and})$ world. Both are semirings, and the Floyd--Warshall update is
really a generic semiring algorithm. Change the two operations and the same
triple loop computes a different fact about the graph. Distances, reachability,
maximum-bottleneck capacity, most-reliable paths -- all the same code, four
different pairs of operators.
\end{intuitionbox}

\begin{algorithm}[H]
\caption{\textsc{Transitive-Closure}($G$)}
\KwIn{Adjacency structure of an $n$-vertex graph $G$.}
\KwOut{Boolean matrix $T$ with $T[i][j] = \texttt{true}$ iff $j$ is reachable from $i$.}
\For{$i \leftarrow 1$ \KwTo $n$}{
  \For{$j \leftarrow 1$ \KwTo $n$}{
    \lIf{$i = j$ \textbf{or} $(i,j) \in E$}{$T[i][j] \leftarrow \texttt{true}$}
    \lElse{$T[i][j] \leftarrow \texttt{false}$}
  }
}
\For{$k \leftarrow 1$ \KwTo $n$}{
  \For{$i \leftarrow 1$ \KwTo $n$}{
    \For{$j \leftarrow 1$ \KwTo $n$}{
      $T[i][j] \leftarrow T[i][j]\ \textbf{or}\ \big(T[i][k]\ \textbf{and}\ T[k][j]\big)$\;
    }
  }
}
\Return $T$\;
\end{algorithm}

Same $\Th(n^3)$ shape, but each operation is a single bit op. If you pack rows
into machine words, the inner loop becomes a bitwise-OR of word-sized chunks,
and you get a real constant-factor speed-up of roughly the word size $w$
(commonly $64$), giving $\Th(n^3 / w)$ in practice -- which is why transitive
closure on bitsets is genuinely fast.

\begin{aestheticbox}
Reachability is the backbone of static analysis. Compilers ask ``can the value
defined here reach this use?''; type checkers and build systems ask ``does
target $A$ transitively depend on target $B$?''; security tools ask ``can
tainted input flow to this dangerous sink?'' All of these are transitive
closures of a graph, and the bitset version of this very loop is what runs
under the hood at scale.
\end{aestheticbox}

\section{Johnson's Algorithm: Fast on Sparse Graphs, Even with Negative Edges}
\label{ch25:johnson}

Floyd--Warshall is a joy on dense graphs, but its $\Th(n^3)$ is wasteful when
the graph is \emph{sparse} -- few edges, $m \ll n^2$. On a sparse graph you
would love to run Dijkstra from each vertex, since one Dijkstra costs
$\Oh(m \log n)$ and $n$ of them cost $\Oh(n m \log n)$, far below $n^3$ when
$m$ is small. There is just one problem: Dijkstra requires \emph{nonnegative}
edge weights, and we want to handle negative ones.

Donald Johnson's idea is to \emph{reweight} the graph so that all
weights become nonnegative, \emph{without changing which paths are shortest},
and then run Dijkstra $n$ times on the reweighted graph. The reweighting itself
is done with one run of Bellman--Ford. Here is how.

\subsection{The reweighting trick}

We want to assign each vertex $v$ a number $h(v)$ -- think of it as an
``elevation'' or potential -- and replace every edge weight by

\[
\hat{w}(u, v) = w(u, v) + h(u) - h(v).
\]

Watch what happens to the weight of an entire path $p = v_0 \to v_1 \to \cdots
\to v_k$. The new weight telescopes:

\[
\hat{w}(p) = \sum_{i=1}^{k}\big(w(v_{i-1}, v_i) + h(v_{i-1}) - h(v_i)\big)
= w(p) + h(v_0) - h(v_k).
\]

Every interior $h$ cancels with its neighbor; only the endpoints survive. So
the new weight of \emph{any} path from $s$ to $t$ is the old weight plus the
same constant $h(s) - h(t)$, regardless of which path you took. This is the
crucial point. Figure~\ref{ch25:fig-telescope} shows the cancellation on a
three-edge path: each $+h(v_i)$ from one edge meets the $-h(v_i)$ of the next.

\begin{figure}[H]
\centering
\begin{tikzpicture}[>=Stealth, node distance=2.0cm]
  \node[circle,draw,fill=sourceblue!18] (v0) {$v_0$};
  \node[circle,draw,right=of v0] (v1) {$v_1$};
  \node[circle,draw,right=of v1] (v2) {$v_2$};
  \node[circle,draw,fill=targetgreen!18,right=of v2] (v3) {$v_3$};
  \draw[->,thick] (v0) -- node[above,font=\scriptsize]{$w_{01}{+}h(v_0){-}h(v_1)$} (v1);
  \draw[->,thick] (v1) -- node[above,font=\scriptsize]{$w_{12}{+}h(v_1){-}h(v_2)$} (v2);
  \draw[->,thick] (v2) -- node[above,font=\scriptsize]{$w_{23}{+}h(v_2){-}h(v_3)$} (v3);
  % cancellation braces
  \node[trapred,font=\scriptsize] at ($(v1)+(0,-0.85)$) {$-h(v_1){+}h(v_1)=0$};
  \node[trapred,font=\scriptsize] at ($(v2)+(0,-0.85)$) {$-h(v_2){+}h(v_2)=0$};
  \node[font=\scriptsize] at ($(v0)!0.5!(v3)+(0,-1.6)$)
    {net reweighted length $=w(p)+h(v_0)-h(v_3)$};
\end{tikzpicture}
\caption{Reweighting telescopes. Each interior potential $h(v_i)$ appears once
with a $+$ and once with a $-$ and cancels; only the endpoint terms $+h(v_0)$
and $-h(v_3)$ survive, so every $v_0$-to-$v_3$ path shifts by the same amount.}
\label{ch25:fig-telescope}
\end{figure}

\begin{focusbox}
Because the additive correction $h(s) - h(t)$ depends only on the endpoints,
\emph{the ranking of paths from $s$ to $t$ is unchanged}. A path that was
shortest in $w$ is still shortest in $\hat{w}$. We have changed all the numbers
but preserved every shortest path. To recover the true distance afterward we
just undo the shift: $\delta(s,t) = \hat\delta(s,t) - h(s) + h(t)$.
\end{focusbox}

\begin{intuitionbox}
Think of $h(v)$ as the height of vertex $v$ and an edge's reweight as
``how much you climb minus how much you descend'' added to its length. Walking
a whole route, all the intermediate ups and downs cancel; only your net
elevation change between start and finish matters. Reweighting is just adding a
conservative ``gravity'' field. Gravity does not change which route between two
fixed points is shortest -- it only re-labels the costs.
\end{intuitionbox}

\subsection{Choosing $h$ so every $\hat{w} \ge 0$}

We need an $h$ that makes $\hat{w}(u,v) = w(u,v) + h(u) - h(v) \ge 0$ for every
edge, i.e. $h(v) - h(u) \le w(u,v)$. That inequality should look familiar: it
is exactly the \emph{triangle / relaxation} condition that shortest-path
distances satisfy. So if we let $h(v)$ be a shortest-path distance from some
fixed source to $v$, the condition holds automatically.

Here is the construction. Add a brand-new vertex $s$ to the graph, with a
zero-weight edge $s \to v$ to every original vertex $v$ (and no edges into $s$).
This new vertex can reach everything and reaches it with finite distance (at
worst $0$). Run Bellman--Ford from $s$. Set $h(v) = \delta(s, v)$, the shortest
distance from $s$ to $v$ in the augmented graph.

Why does this satisfy $h(v) - h(u) \le w(u, v)$? Because shortest distances
obey the triangle inequality: the best way to reach $v$ is no worse than the
best way to reach $u$ followed by the edge $(u,v)$, so
$\delta(s,v) \le \delta(s,u) + w(u,v)$, which rearranges to exactly what we
need. And Bellman--Ford from $s$ is also where we \emph{detect a negative
cycle}: if it reports one, we stop and announce that all-pairs shortest paths
are undefined.

\begin{figure}[H]
\centering
\begin{tikzpicture}[>=Stealth,node distance=1.6cm]
  \node[circle,draw,fill=quotientorange!22] (s) {$s$};
  \node[circle,draw,fill=blue!12,above right=0.4cm and 2cm of s] (a) {$1$};
  \node[circle,draw,fill=blue!12,right=2.6cm of s] (b) {$2$};
  \node[circle,draw,fill=blue!12,below right=0.4cm and 2cm of s] (c) {$3$};
  \draw[->,dashed,gray] (s) -- node[above,sloped,font=\scriptsize]{0} (a);
  \draw[->,dashed,gray] (s) -- node[above,font=\scriptsize]{0} (b);
  \draw[->,dashed,gray] (s) -- node[below,sloped,font=\scriptsize]{0} (c);
  \draw[->,thick] (a) -- node[right,font=\scriptsize]{$-4$} (c);
  \draw[->,thick] (b) -- node[above,font=\scriptsize]{$2$} (a);
  \draw[->,thick] (c) to[bend left=20] node[below,font=\scriptsize]{$5$} (b);
\end{tikzpicture}
\caption{The augmented graph $G'$ for Johnson's algorithm. A synthetic source
$s$ (orange) sends a dashed $0$-weight edge to every real vertex; one
Bellman--Ford run from $s$ then yields a potential $h(v)=\delta(s,v)$ for each
vertex in a single sweep.}
\label{ch25:fig-johnson-source}
\end{figure}

The dashed $0$-weight edges from the synthetic source $s$ let Bellman--Ford
compute one $h$ value per vertex in a single pass over the whole graph
(Figure~\ref{ch25:fig-johnson-source}).

\subsection{The full algorithm}

\begin{algorithm}[H]
\caption{\textsc{Johnson}($G$, $w$)}
\KwIn{Directed weighted graph $G = (V, E)$ with weights $w$, possibly negative.}
\KwOut{Distance matrix $D$, or a report that a negative-weight cycle exists.}
form $G'$ by adding a new vertex $s$ and edges $(s, v)$ of weight $0$ for all $v \in V$\;
\If{\textsc{Bellman-Ford}($G'$, $w$, $s$) finds a negative-weight cycle}{
  \Return ``negative-weight cycle''\;
}
\ForEach{$v \in V$}{$h(v) \leftarrow \delta(s, v)$ \tcp*{from Bellman--Ford}}
\ForEach{edge $(u, v) \in E$}{$\hat{w}(u, v) \leftarrow w(u, v) + h(u) - h(v)$}
\ForEach{$u \in V$}{
  run \textsc{Dijkstra}($G$, $\hat{w}$, $u$) to get $\hat\delta(u, v)$ for all $v$\;
  \ForEach{$v \in V$}{
    $D[u][v] \leftarrow \hat\delta(u, v) - h(u) + h(v)$\;
  }
}
\Return $D$\;
\end{algorithm}

\subsection{A worked example of the reweighting}

\begin{labbox}
Use the three-vertex graph from the diagram: edges $2 \to 1$ (weight $2$),
$1 \to 3$ (weight $-4$), $3 \to 2$ (weight $5$). Add source $s$ with
$0$-weight edges to $1, 2, 3$ and run Bellman--Ford. The cheapest way to reach
each vertex from $s$:
\[
h(1) = \delta(s,1) = 0,\quad h(2) = \delta(s,2) = 0,\quad h(3) = \delta(s,3) = -4.
\]
(Vertex $3$ is reached via $s \to 1 \to 3 = 0 + (-4) = -4$; everything else is
reached directly at cost $0$, since the synthetic edges already cost $0$.)
Now reweight each real edge with $\hat{w}(u,v) = w(u,v) + h(u) - h(v)$:
\[
\hat{w}(2,1) = 2 + 0 - 0 = 2,\quad
\hat{w}(1,3) = -4 + 0 - (-4) = 0,\quad
\hat{w}(3,2) = 5 + (-4) - 0 = 1.
\]
Every reweighted edge is nonnegative -- exactly as promised -- so Dijkstra is
now legal. And the $-4$ edge, which would have broken Dijkstra, became a
harmless $0$. To read a true distance back out, say from $1$ to $2$: Dijkstra
on $\hat{w}$ gives $\hat\delta(1,2) = \hat{w}(1,3) + \hat{w}(3,2) = 0 + 1 = 1$;
then $\delta(1,2) = \hat\delta(1,2) - h(1) + h(2) = 1 - 0 + 0 = 1$, while the
true path $1 \to 3 \to 2$ costs $-4 + 5 = 1$. They match.
\end{labbox}

\subsection{Why it is correct, and what it costs}

\begin{formalbox}
\textbf{Claim.} With $h(v) = \delta(s, v)$, every reweighted edge satisfies
$\hat{w}(u,v) \ge 0$, and for every pair $(u, v)$,
$\delta(u, v) = \hat\delta(u, v) - h(u) + h(v)$.

\emph{Nonnegativity.} The triangle inequality gives $\delta(s,v) \le
\delta(s,u) + w(u,v)$, i.e. $h(v) - h(u) \le w(u,v)$, i.e.
$\hat{w}(u,v) = w(u,v) + h(u) - h(v) \ge 0$.

\emph{Distance recovery.} For any path $p$ from $u$ to $v$, telescoping gives
$\hat{w}(p) = w(p) + h(u) - h(v)$. The additive term is the same for all $p$,
so the $\hat{w}$-shortest path is a $w$-shortest path; taking the minimum over
$p$ on both sides yields $\hat\delta(u,v) = \delta(u,v) + h(u) - h(v)$, which
rearranges to the stated recovery formula.
\end{formalbox}

\begin{complexitybox}
\textbf{Johnson's algorithm.}
\begin{itemize}
\item One Bellman--Ford from $s$: $\Oh(n m)$.
\item Reweighting all edges: $\Oh(m)$.
\item $n$ runs of Dijkstra with a binary heap: $\Oh\big(n\,(m \log n)\big) =
\Oh(n m \log n)$; with a Fibonacci heap, $\Oh\big(n\,(m + n \log n)\big) =
\Oh(n m + n^2 \log n)$.
\item Total: $\Oh(n m \log n)$ (binary heap), dominated by the Dijkstra runs.
\item Space $\Oh(n^2)$ for the output, $\Oh(n + m)$ working.
\item Handles negative edges; detects negative cycles in the Bellman--Ford
phase.
\end{itemize}
\end{complexitybox}

\begin{intuitionbox}
The division of labor is the whole point. Bellman--Ford is slow but robust: it
tolerates negative edges, so we use it \emph{once} to neutralize them by
reweighting. Dijkstra is fast but fragile: it cannot tolerate negative edges,
so we run it $n$ times only \emph{after} the graph has been made nonnegative.
Johnson's algorithm is a marriage that takes the best trait from each partner.
\end{intuitionbox}

\section{Choosing an Algorithm}
\label{ch25:choosing}

You now have several tools. Which one should you reach for?

\begin{center}
\begin{tabular}{l l l}
\toprule
Method & Time & Best when \\
\midrule
Repeated Dijkstra & $\Oh(n m \log n)$ & sparse, weights $\ge 0$ \\
Johnson & $\Oh(n m \log n)$ & sparse, some negative edges \\
Floyd--Warshall & $\Th(n^3)$ & dense, any weights, want simplicity \\
Matrix (squaring) & $\Th(n^3 \log n)$ & mainly of theoretical interest \\
Transitive closure & $\Th(n^3 / w)$ & only reachability needed \\
\bottomrule
\end{tabular}
\end{center}

The rough rule: \textbf{dense graph $\Rightarrow$ Floyd--Warshall};
\textbf{sparse graph $\Rightarrow$ Johnson} (or repeated Dijkstra if all weights
are already nonnegative). The crossover is around $m \approx n^2 / \log n$.
Below that density, $n m \log n < n^3$ and Johnson wins; above it,
Floyd--Warshall's tiny constant factor and cache-friendly triple loop make it
the practical champion.

\begin{historybox}
This material is CLRS Chapter 25, ``All-Pairs Shortest Paths,'' which builds
directly on CLRS Chapter 24 (single-source) and the graph foundations of
Chapter 22. The matrix-multiplication formulation and Floyd--Warshall are CLRS
\S 25.1--25.2; the transitive-closure variant is in \S 25.2; Johnson's
algorithm is \S 25.3. Bernard Roy (1959), Stephen Warshall (1962), and Robert
Floyd (1962) arrived at the triple-loop algorithm independently -- Warshall for
reachability, Floyd for distances -- which is why you will sometimes see it
called the Roy--Floyd--Warshall algorithm. Donald Johnson published his
reweighting algorithm in 1977. The deeper observation, that all of these are
the \emph{same} algorithm over different semirings, is part of the algebraic
path problem and connects onward to Kleene's theorem in automata theory.
\end{historybox}

\begin{aestheticbox}
All-pairs shortest paths quietly powers things you use daily. Route planners
precompute distance tables over regions so a query is a lookup, not a fresh
search. The classic distance-vector routing protocols (RIP and friends) are
distributed cousins of Bellman--Ford run continuously across the internet.
Network-analysis tools compute \emph{betweenness centrality} -- how often a node
sits on shortest paths between others -- which needs all-pairs distances and is
used to find influential people, fragile links, and bottleneck routers. And in
program analysis, the transitive-closure form decides the reachability questions
that compilers and security scanners depend on.
\end{aestheticbox}

\section{Summary}
\label{ch25:summary}

\begin{recoverybox}
The through-line of the chapter, in one place.
\begin{itemize}
\item The problem: a full $n \times n$ table of shortest distances between all
ordered pairs, plus optional predecessors for path reconstruction.
\item Organize paths by \emph{number of edges} and you get a min-plus matrix
product; squaring gives $\Th(n^3 \log n)$.
\item Organize paths by \emph{which intermediate vertices are allowed} and you
get Floyd--Warshall, a clean in-place $\Th(n^3)$ triple loop.
\item Swap $(\min, +)$ for $(\textsc{or}, \textsc{and})$ and the same loop
computes transitive closure (reachability) on bits.
\item For sparse graphs with negative edges, Johnson reweights once with
Bellman--Ford to kill negativity, then runs Dijkstra $n$ times for
$\Oh(n m \log n)$.
\item Negative cycles make distances $-\infty$; watch the Floyd--Warshall
diagonal, or trust Bellman--Ford in Johnson, to detect them.
\end{itemize}
\end{recoverybox}

\section*{Exercises}
\addcontentsline{toc}{section}{Exercises}

\begin{exercise}
Run \textsc{Faster-APSP} (repeated squaring) by hand on the four-vertex graph
from the Floyd--Warshall lab. Write $L^{(1)} = W$, then $L^{(2)}$, then
$L^{(4)}$. Confirm that $L^{(4)}$ equals the final matrix $D^{(4)}$ we obtained
with Floyd--Warshall. Why is computing $L^{(4)}$ (rather than stopping at
$L^{(3)}$) the right call for $n = 4$?
\end{exercise}

\begin{exercise}
Carry the predecessor matrix $\Pi$ through the full Floyd--Warshall trace in the
lab. Then use it to print the actual shortest path from vertex $2$ to vertex
$4$. (You should recover $2 \to 3 \to 4$ with total weight $3$.)
\end{exercise}

\begin{exercise}
Explain precisely why the loop order in Floyd--Warshall must have $k$ outermost.
Construct a tiny graph (three vertices is enough) on which putting the $k$ loop
innermost produces a wrong answer, and show the wrong number it computes.
\end{exercise}

\begin{exercise}
The in-place version of Floyd--Warshall reuses one matrix. We justified this by
arguing the $k$-th row and $k$-th column do not change during pass $k$. Prove
that claim carefully from the recurrence
$d^{(k)}_{ij} = \min(d^{(k-1)}_{ij}, d^{(k-1)}_{ik} + d^{(k-1)}_{kj})$ by
specializing to $j = k$ and to $i = k$.
\end{exercise}

\begin{exercise}
Modify Floyd--Warshall so that instead of \emph{shortest} paths it computes, for
every pair $(i, j)$, the maximum-bottleneck path: the path that maximizes the
minimum edge weight along it (useful for finding the widest channel between two
network nodes). What two operations replace $\min$ and $+$? What is the base
case? This is another semiring -- name the operations.
\end{exercise}

\begin{exercise}
A negative-weight cycle leaves a fingerprint on the Floyd--Warshall diagonal.
Run Floyd--Warshall on the two-vertex graph with edges $1 \to 2$ (weight $1$)
and $2 \to 1$ (weight $-3$). Show that after the algorithm finishes, $D[1][1]$
and $D[2][2]$ are negative, and explain how you would scan for this.
\end{exercise}

\begin{exercise}
In Johnson's algorithm we added a synthetic source $s$ with $0$-weight edges to
every vertex. Suppose instead you used edges of weight $+5$. Would the algorithm
still be correct? What about weight $-5$? Argue from the role $h(v) =
\delta(s,v)$ plays, and state exactly which property of $h$ you must preserve.
\end{exercise}

\begin{exercise}
Work out the density crossover between Johnson and Floyd--Warshall more
carefully. Setting $n m \log n = n^3$, solve for $m$ in terms of $n$. For
$n = 1000$, roughly how many edges does a graph need before Floyd--Warshall
becomes the faster choice? (Use $\log_2 1000 \approx 10$.)
\end{exercise}

\begin{exercise}
(\emph{Challenge.}) Show how to use the predecessor information from a single
run of Floyd--Warshall to detect, for a graph \emph{with} a negative cycle, an
actual vertex sequence forming a negative cycle -- not just the fact that one
exists. Hint: a negative diagonal entry $D[i][i] < 0$ means $i$ is on such a
cycle; follow predecessors from $i$ back to $i$.
\end{exercise}

\begin{exercise}
(\emph{Challenge.}) The min-plus product has no known truly sub-cubic
\emph{combinatorial} algorithm, but there is a clever $\Oh(n^3 / \log n)$-style
speed-up and connections to the famous APSP-vs-min-plus equivalence. Without
looking it up, explain why Strassen's $\Oh(n^{2.81})$ matrix multiplication
cannot simply be dropped in to speed up min-plus APSP. (Hint: what algebraic
operation does Strassen rely on that $\min$ does not have?)
\end{exercise}
